\paragraph{Definice (stručně)} 
Autenticita: úloha v reálném kontextu měřící schopnost aplikovat znalosti v praktických situacích.

\paragraph{Proč přehodnotit zadání směrem k autenticity}
Autentická zadání motivují studenty a snižují možnost „vyřešení“ jen pomocí generativní AI, protože úspěch závisí na kontextu, procesu a komunikaci výsledků.

\paragraph{Postup redesignu úlohy (praktický návod)}
\begin{enumerate}
  \item Vymezte reálný kontext a adresáta; navrhněte artefakty a milníky (drafty, protokoly, reporty).
  \item Zahrňte rubriku hodnotící proces, relevanci, komunikaci výsledků a reflexi o použití AI (deklarace).
  \item Uskutečnějte pilot v jednom kurzu, ověřte fungování a zapracujte lessons‑learned; při plánování pilotu lze využít Copilot‑typ agenty — vždy s lidskou kontrolou \cite{kb_ai_101_tipu_pdf_04e4a115_page_48_1}.
\end{enumerate}

\paragraph{Krátký příklad} Zadání: analyzujte spotřebu energie labu, navrhněte dvě opatření; artefakty: metodika (milník 1), předběžná analýza (milník 2), finální report + prezentace + reflexe o AI; rubrika: relevance (30\%), kvalita analýzy (30\%), komunikace (20\%), reflexe/AI (20\%).

\paragraph{Rychlý checklist pro vyučující}
\begin{itemize}
  \item Je úloha vázána na adresáta? \item Jsou definovány milníky a artefakty? \item Obsahuje rubrika procesu a reflexi AI? \item Byl naplánován pilot a mechanismus spot‑checků? Pokud ano, zvažte nástroje pro plánování pilotu \cite{kb_ai_101_tipu_pdf_04e4a115_page_48_1}.
\end{itemize}